% Slide — Próximos passos
\begin{frame}{Onde isso está sendo validado agora}
  \small As 18 políticas, com a função de aptidão em seu estágio mais
  recente (Eq.~4.2/Zabraoui $+$ TR/BE), estão sendo reexecutadas em
  \textbf{três bases}, com o mesmo critério de otimização em todas:
  \vspace{0.4em}

  \tabstyle\footnotesize
  \begin{tabular}{@{}lll@{}}
    \toprule
    \textbf{Base} & \textbf{Escopo} & \textbf{Status} \\
    \midrule
    Bahia (Experimento 2, oficial) & 145 séries & concluída \\
    ``bot'' (interna completa, 27 estados) & 4.869 séries & em andamento \\
    M5 (Walmart, benchmark externo, sem cortes) & 30.490 séries & em andamento \\
    \bottomrule
  \end{tabular}
  \vspace{1em}

  \begin{block}{Por que reexecutar a Bahia}
    \small A extensão TR/BE (estágio 3 da aptidão) foi aplicada
    \textbf{durante} a execução anterior da Bahia -- deixando etapas já
    concluídas (clássica, meta-heurística, híbrida) com a aptidão antiga.
    Reexecutada do zero para que as \textbf{três bases otimizem pelo mesmo
    critério}, condição necessária para qualquer comparação entre elas.
  \end{block}

  \note{
    Este slide fecha o círculo entre a implementação apresentada e o
    trabalho em andamento. Como a extensão de aptidão com TR e BE foi a
    mudança mais recente, aplicada no meio de uma execução de produção da
    Bahia, essa execução específica ficou com etapas inconsistentes entre
    si -- as que já tinham rodado usaram o critério antigo. A decisão foi
    reexecutar a Bahia do zero com o critério novo do início ao fim, para
    garantir que qualquer comparação futura entre Bahia, "bot" e M5 esteja
    otimizando exatamente o mesmo objetivo nas três. Isso não é
    apresentado como resultado -- é o estado do trabalho de validação no
    momento desta apresentação.
  }
\end{frame}
